% ============================================================================================
% J.-P. Benzécri, « Sur les variétés localement affines et localement projectives »,
% Bulletin de la S. M. F. 88 (1960), p. 229-332 (thèse Sc. math., Paris, 1960).
% DOI 10.24033/bsmf.1551 ; Numdam BSMF_1960__88__229_0.
%
% TRANSCRIPTION, not the author's source. Transcribed on 2026-11-28 by a Claude session
% (T-bz-chapter3, DEC-195) from the Numdam scan papers/benzecri/sources/bsmf-1960-88-229/benzecri1960.pdf
% (sha256 77e0b535382fa33109b7615fb9fa336151a86723ab8210b92b91c01b0e306306), reading the page
% images, with the scan's OCR text layer as a help only. TO BE CHECKED BY THE AUTHOR.
%
% Scope: Chapter 3 « Compléments sur les convexes » complete, journal pp. 259-283
% (PDF page = journal page - 227; PDF pp. 32-56): from its title (p. 259, where Chapter 2 ends)
% to the foot of p. 283; the title of Chapter 4 opens p. 284. Sections 1 to 6.
%
% Conventions of the transcription (those of benzecri1960_ch1.tex and benzecri1960_ch2.tex).
%  * The French text is the printed text, misprints included; each misprint we noticed is kept and
%    flagged by a comment `% [sic] ...` giving the evident intended reading; `% [sic?]` marks a
%    reading that is probable but not certain. Words cut at a line end are rejoined. Nothing
%    illegible was met; `[illisible]` would mark it.
%  * `% --- p. NNN ---` marks the start of journal page NNN.
%  * Each numbered statement is in its environment, the printed number as the optional argument
%    (with its prime: `1$'$`, `1$''$`). Where a statement ends is printed only for the THÉORÈMES,
%    set in italics; every other statement of Chapter 3 is set in roman, so the end of its
%    environment is an editorial choice, made at the first paragraph that is commentary or
%    notation rather than part of the statement. The fonts the environments print are the
%    preamble's, not the scan's. Proofs (printed without a heading, usually « En effet, ... »,
%    or headed « DÉMONSTRATION. — », kept as text) are in `proof` environments, idem.
%  * Unnumbered DÉFINITION, LEMME, COROLLAIRE and REMARQUE blocks are environments without
%    argument; the numbered COROLLAIRE 1 and 2 of § 5 carry their number. No new environment was
%    needed; `defin`, `lemme` and `corollaire` print no stray space when they have no argument
%    (\bznum), and `corollaire` accepts a number.
%  * Where a proof is printed away from its statement (§ 2, Prop. 3, proved after Déf. 4;
%    § 2, Theorem 1, whose proof is interrupted by the LEMME and its proof), the `proof`
%    environment carries its target as an optional argument « Proof of » + \ref, editorial and not
%    printed, as in Chapter 2. These are the only \ref outside a \bzcf.
%  * \bzcf{printed text}{\ref{...}} marks an explicit cross-reference the source makes to a
%    statement: it prints the printed text only; the \ref names the statement the PRINTED text
%    designates, even where that is a misprint (flagged [sic]). Chapter 1 statements are named by
%    their labels in legacy/labeled/benzecri1960_ch1.tex; Chapter 3 statements by the labels of
%    the scheme bz<kind>:c3s<section><initial><number>_<slug> listed in NOTES.md (the \label's
%    themselves are not in this file). A reference of a proof to its own statement, to the
%    hypotheses a, b, c, d of § 3 (free text), to Bourbaki, or to Chapter 4 is not marked.
%  * $\R$ is the printed bold R (the reals); $\mathring{H}$ the printed H with a ring (an open
%    hemisphere; likewise $\mathring{U}$, $\mathring{V}$, $\mathring{B}$); \intr{X} the printed
%    ring over an arc spanning X (the interior of X); $\complement$ the printed complement sign;
%    $\Dc$ the printed script D of § 4; $w$, $v$, $v_i$ the printed script w and v (classes of
%    § 2, convexes of S(E), points of § 6); $)a, b($ the printed open interval, kept as printed.
% ============================================================================================
\documentclass[11pt]{article}
\usepackage[utf8]{inputenc}
\usepackage[T1]{fontenc}
\usepackage[french]{babel}
\usepackage{amsmath,amssymb}
% Layout only (two lines of Chapter 3 would otherwise be overfull); no effect on the text.
\setlength{\emergencystretch}{2em}
\newcommand{\R}{\boldsymbol{R}}
\newcommand{\Rc}{\mathcal{R}}
\newcommand{\bzcf}[2]{#1}
% \intr{X} is the printed ring over an arc spanning X: the interior of X (§ 3, Props. 9, 10).
% \Dc is the printed script D (the correspondence of facettes of § 4, Props. 6-9).
\newcommand{\Dc}{\mathcal{D}}
\newcommand{\intr}[1]{\overset{\circ}{\widehat{#1}}}
% \bznum{n} prints " n" and nothing for an empty argument (unnumbered THÉORÈME, REMARQUE).
\newcommand{\bznum}[1]{\if\relax\detokenize{#1}\relax\else\ #1\fi}
\newenvironment{resume}[1][]{\par\medskip\noindent #1. }{\par\medskip}
\newenvironment{defin}[1][]{\par\medskip\noindent\textsc{Définition\bznum{#1}.} --- }{\par\medskip}
\newenvironment{prop}[1][]{\par\medskip\noindent\textsc{Proposition #1.} --- \itshape}{\par\medskip}
\newenvironment{theoreme}[1][]{\par\medskip\noindent\textsc{Théorème\bznum{#1}.} --- \itshape}{\par\medskip}
\newenvironment{lemme}[1][]{\par\medskip\noindent\textsc{Lemme\bznum{#1}.} --- \itshape}{\par\medskip}
\newenvironment{corollaire}[1][]{\par\medskip\noindent\textsc{Corollaire\bznum{#1}.} --- \itshape}{\par\medskip}
\newenvironment{remarque}[1][]{\par\medskip\noindent\textsc{Remarque\bznum{#1}.} --- }{\par\medskip}
\newenvironment{exemple}[1][Exemple]{\par\medskip\noindent\textsc{#1.} --- }{\par\medskip}
\newenvironment{enonce}{\par\medskip\noindent\itshape}{\par\medskip}
% The optional argument of `proof` is accepted and not printed (room for a ledger target).
\newenvironment{proof}[1][]{\par}{\par}

\begin{document}

% --- p. 259 ---
% (Chapter 2 ends at the top of p. 259; see benzecri1960_ch2.tex.)
\section*{Chapitre 3. Compléments sur les convexes.}

Ce chapitre contient la démonstration de propriétés des convexes affins et projectifs dont il sera
fait usage au chapitre 5 dans l'étude des convexes balayables.

\subsection*{1. Facettes et variétés d'appui.}

Dans ce paragraphe, les notions de facettes et de variétés d'appui sont définies pour les convexes
d'un espace affin et de la sphère projective associée.

\begin{defin}[1]
Sur un ensemble convexe $W$, dans un espace vectoriel, on définit la relation de préordre
$\alpha(W)$ : $A\,\alpha(W)\,B$ s'il existe un segment $(A', B)\subset W$ tel que
$A\in)A', B($, ou si $A=B$.
\end{defin}

On peut simplifier l'énoncé de la définition 1 en faisant la convention suivante : on appelle
intérieur d'un segment $AB$, l'intérieur de ce segment, muni de la topologie induite par la
variété linéaire support de l'ensemble $\{A, B\}$; cet intérieur se note $)A, B($ ou $)B, A($; si
$A\neq B$, on retrouve la définition habituelle de l'intervalle ouvert $)A, B($; si $A=B$, le
support de $\{A\}$ n'est autre que la variété à $0$ dimension $\{A\}$, et l'on a
\[ )A, A(=(A, A)=\{A\}. \]
L'énoncé de la définition 1 devient :

$A\,\alpha(W)\,B$ s'il existe un segment $(A', B)\subset W$ tel que $A\in)A', B($.

Pour $A=B$ on pose $A'=A$ et l'on a $A\in)A, A($.

Dans ce qui suit, nous écrirons toujours $\alpha$ pour $\alpha(W)$ si aucune confusion n'en peut
résulter.

% The verification that α is a preorder is printed after the definition, without a heading;
% it is put in a `proof` environment, which the ledger attaches to Déf. 1 (the nearest statement).
\begin{proof}
Pour montrer qu'on a bien une relation de préordre, il faut démontrer
% --- p. 260 ---
que $A\,\alpha\,B$ et $B\,\alpha\,C$ impliquent $A\,\alpha\,C$. En effet, si l'on a dans $W$,
$A'$ et $B'$ tels que
\[ \begin{aligned} &A\in)A', B(\\ &B\in)B', C( \end{aligned} \]
$A'$, $B'$, $C$ engendrent un convexe : $A$ est intérieur à ce convexe supposé muni de la
topologie de la variété linéaire support, et le support de $AC$ le coupe suivant un segment
auquel $A$ est intérieur.
\end{proof}

\begin{defin}[2]
La relation d'équivalence associée à la relation de préordre $\alpha(W)$ se note $F_W(A, B)$ :
les classes d'équivalence correspondantes sont appelées \textit{$W$-facettes} ($W$ ne sera noté
que si une confusion est à craindre).

La relation de préordre $\alpha$ définit canoniquement sur l'ensemble des facettes une relation
d'ordre que nous noterons aussi $\alpha$.
\end{defin}

\begin{remarque}
On peut donner de la définition de $F_W(A, B)$ l'énoncé suivant : $n$ points $A_1$, \ldots, $A_n$
de $W$ appartiennent à une même facette si, pour la topologie de la variété linéaire $L$ support
de $\{A_1, \ldots, A_n\}$, $A_1$, \ldots, $A_n$ sont intérieurs à $L\cap W$.
\end{remarque}
% « A_1, .... A_n » is printed with a period after the dots (read « A_1, ..., A_n »).

\begin{prop}[1]
Soit $F$ une facette d'un convexe $W$, et soit $L$ la variété linéaire support de $F$; $F$ est un
ouvert de $L$, supposé muni de sa topologie localement convexe la plus fine.
\end{prop}

\begin{proof}
Il suffit de démontrer que si $A\in F$, et si $l$ est une sous-variété linéaire de dimension finie
telle que $A\in l$ et $l\subset L$, alors $A$ est intérieur à $l\cap F$ pour la topologie de $l$.
Or, cela résulte de \bzcf{la remarque ci-dessus}{\ref{bzrem:c3s1r1_facette_points_interieurs}}.
\end{proof}

\begin{prop}[2]
Soient $F$ une facette de $W$, $L$ son support, et soient $A\in F$, $B\in W$; alors
$A\,\alpha\,B$ si et seulement si $B\in L$.
\end{prop}

\begin{proof}
En effet, si $B\in L$, la variété linéaire engendrée par $A$ et $B$ coupe $W$ suivant une partie
convexe, qui contient $B$ et à laquelle $A$ est intérieur, donc $A\,\alpha\,B$; si
$A\,\alpha\,B$, la variété linéaire engendrée par $A$ et $B$ est coupée par $F$, donc par $L$
suivant une partie convexe d'intérieur non vide, donc $B\in L$.
\end{proof}

\begin{prop}[3]
Soit $F$, $F'$ deux facettes de $W$ de supports $L$ et $L'$; on a $F\,\alpha\,F'$ si et seulement
si $L'\subset L$; on a $L=L'$ si et seulement si $F=F'$.
\end{prop}
% (No proof of Prop. 3 is printed.)

\begin{prop}[4]
L'ensemble des facettes muni de la relation d'ordre $\alpha$ est filtrant à gauche.
\end{prop}

\begin{proof}
En effet, soient $F$ et $F'$ deux facettes $A\in F$, $A'\in F'$, $B$ un point intérieur au
segment $AA'$ dans son support : $B\,\alpha\,A$; $B\,\alpha\,A'$. La facette $F''$ de $B$ est donc
telle que $F''\,\alpha\,F'$; $F''\,\alpha\,F$.
\end{proof}

\begin{corollaire}
L'ensemble des variétés linéaires supports de facettes du convexe $W$ est filtrant à droite pour
la relation d'inclusion.
\end{corollaire}

% --- p. 261 ---
\begin{defin}[3]
Une variété linéaire $S$ est dite \textit{variété d'appui} de $W$ si $S\cap W$ est non vide et
saturé pour la relation d'équivalence $F_W$ (cette définition équivaut à celle de
\textsc{Bourbaki}, \textit{Espaces vectoriels topologiques}, chap. II, p. 83).
\end{defin}

D'après \bzcf{la proposition 3}{\ref{bzprop:c3s1p3_ordre_supports}}, il est équivalent de dire
que $A\in S\cap W$ et $A\,\alpha\,B$ impliquent $B\in S\cap W$.

L'espace tout entier est une variété d'appui, le support d'une facette est une variété d'appui;
la réciproque n'est pas vraie : on trouve facilement, dans un espace de dimension finie, un
convexe $W$ et une variété d'appui $S$ telle que $S\cap W$ ait un support distinct de $S$; nous
allons donner un exemple en dimension infinie d'un convexe $W$ et d'une variété d'appui $S$ de $W$
telle que $S\cap W$ engendre $S$, mais qu'il n'existe aucune facette dont le support soit $S$.
Considérons dans un espace vectoriel de dimension $>\aleph_0$ un système dénombrable de vecteurs
linéairement indépendants $OA_i$ issus de l'origine $O$ : l'enveloppe convexe de $O$ et de $A_i$
est un $\aleph_0$ simplexe $S$ dont le support est de dimension infinie : chaque facette cependant
est de dimension finie; tout point $M$ de $S$ est en effet barycentre d'un nombre fini de points
de l'ensemble $\{O, A, \ldots\}$ affectés de masses non nulles dont la somme est $1$; ces points
en nombre fini engendrent le support de la facette de $M$. En dimension finie. on a toutefois la
% [sic] « En dimension finie. on a » : read « En dimension finie, on a » (period for comma).
% « {O, A, ...} » as printed (read {O, A_1, ...}); « S » names both the variété d'appui and the
% ℵ0-simplex, as printed.

\begin{prop}[5]
Si $W$ est un convexe dont le support $E$ est de dimension finie, pour toute variété d'appui $L$ à
$W$ il existe une facette $F$ dont le support $S$ est tel que $S\cap W=L\cap W$.
\end{prop}

\begin{proof}
Désignons par $S$ le support de $L\cap W$; alors $S\cap W=L\cap W\neq\emptyset$; de plus,
$S\cap W$ est un convexe de support $S$ qui, pour la topologie de l'espace $S$ de dimension finie,
admet donc au moins un point interne dont la facette $F$ aura $S$ pour support.
\end{proof}

\begin{remarque}
1\textsuperscript{o} $E$ est le support d'une facette de $W$ qui n'est autre que son intérieur;

2\textsuperscript{o} Si $W$ est un convexe compact de dimension finie, $L\cap W$ n'est autre que
la fermeture compacte $E$ de $F$.
% [sic] « la fermeture compacte E de F » : read « F̄ » (the closure of F); E is the support of W.
\end{remarque}

\begin{prop}[6]
Soit $L$ une variété linéaire, $W$ un convexe
\[ W'=W\cap L=\emptyset. \]
% [sic] « W' = W ∩ L = ∅ » : read « ≠ ∅ » (the proposition is about the convex W' = W ∩ L).
Alors la relation $\alpha(W')$ est induite par la relation $\alpha(W)$.
\end{prop}

\begin{proof}
Cela résulte immédiatement de la définition de $\alpha$. Des \bzcf{propositions
3}{\ref{bzprop:c3s1p3_ordre_supports}} et 6 il résulte que, si $L$ est une variété d'appui à $W$,
la facette d'un point $M\in L\cap W$, relativement au convexe $W$ n'est autre que sa facette
relativement au convexe $L\cap W$.
\end{proof}

Nous allons maintenant étudier le cas des convexes de la sphère.

% --- p. 262 ---
\begin{defin}[4]
Soit $k$ une partie convexe de $S(E)$; $k'=F(k)+\{O\}$ est un convexe de $E$; $\alpha(k')$
(\bzcf{\textit{cf.} \S\,1, déf. 1}{\ref{bzdef:c3s1d1_preordre_alpha}}) induit sur $F(k)$ une
relation de préordre compatible avec la projection sur $k$. Nous noterons $\alpha(k)$ la relation
de préordre ainsi définie sur $k$.

La relation d'équivalence associée à la relation $\alpha(k)$ se note $F(\alpha(k))$; les classes
d'équivalence correspondantes sont appelées \textit{$k$-facettes} et la relation de préordre
$\alpha$ définit canoniquement sur l'ensemble des facettes une relation d'ordre que nous noterons
aussi $\alpha$.

La relation $\alpha(k)$ peut être définie directement :
\[ a\,\alpha(k)\,b \]
s'il existe un segment $(a', b)\subset k$ tel que $a\in)a', b($ ou si $a$ et $b$ sont
identiques ou diamétralement opposés.
\end{defin}
% « § 1, déf. 1 » is the present § 1 (Chapter 3), where α(W) is defined.

\begin{defin}[5]
Soit $L$ le sous-espace vectoriel facette de $O$ dans le convexe $k'=F(k)+\{O\}$; la projection
sur $k$ de $L-\{O\}$ est une sous-variété linéaire, éventuellement vide, ensemble de tous les
points de $S(E)$ qui appartiennent à $k$ avec leur point diamétralement opposé. Cette
sous-variété est appelée le \textit{noyau} de $k$. Un convexe strict est un convexe de noyau vide
(\bzcf{\textit{cf.} \S\,1, déf. 9}{\ref{bzdef:c1s1d9_convexe_strict}}).
\end{defin}
% « § 1, déf. 9 » : the present § 1 has no Déf. 9; the convexe strict is Chapter 1, § 1, déf. 9,
% which the \ref names.

Les \bzcf{propositions 2}{\ref{bzprop:c3s1p2_alpha_support}},
\bzcf{3}{\ref{bzprop:c3s1p3_ordre_supports}}, \bzcf{4}{\ref{bzprop:c3s1p4_facettes_filtrant}}
et \bzcf{le corollaire}{\ref{bzcor:c3s1p4c_supports_filtrant}} sont valables pour les convexes de
$S(E)$. \bzcf{La proposition 1}{\ref{bzprop:c3s1p1_facette_ouverte}} aussi, si l'on convient
d'appeler topologie localement convexe la plus fine sous-variété $l$ de $S(E)$ la topologie
quotient de celle induite sur $F(l)$ par la topologie localement convexe la plus fine de
$F(l)+\{O\}$.
% « la plus fine sous-variété l » as printed (read « la plus fine de la sous-variété l »).

\begin{defin}[6]
Une variété linéaire $l$ est dite \textit{variété d'appui pour $k$} s'il existe un sous-espace
vectoriel $L$ de $E$, variété d'appui de $F(k)+\{O\}$, et
\[ l=F^{-1}(L-\{O\}). \]
\end{defin}

En particulier, le noyau de $k$ est une variété d'appui de $k$. Si ce noyau est vide, il
constitue une variété d'appui, ou encore une facette, de $k$ de dimension $(-1)$.

Les \bzcf{propositions 5}{\ref{bzprop:c3s1p5_appui_facette_dim_finie}} et
\bzcf{6}{\ref{bzprop:c3s1p6_alpha_induit}} sont également valables pour les convexes de $S(E)$.

\subsection*{2. Sommes de convexes.}

Comme au précédent paragraphe, on étudie le cas affin, puis le cas projectif.

\begin{defin}[1]
Soient $E$ un espace vectoriel, $W$ un convexe de $E$, $V_i$ ($i\in I$) une famille finie de
convexes de $E$, $p$ l'application de $\prod_{i\in I}V_i$ dans $E$ défini par
\[ \{x_i\}_{i\in I}\in\prod_{i\in I}V_i,\qquad p(\{x_i\}_{i\in I})=\sum_{i\in I}x_i, \]
% --- p. 263 ---
$W$ est dit \textit{somme des convexes $V_i$} si $p$ est un isomorphisme de $\prod_{i\in I}V_i$
sur $W$ (pour les structures affines). Nous noterons
\[ W=\sum_{i\in I}V_i\qquad\text{ou}\qquad W=V_1+V_2+\ldots+V_n. \]
Si $x\in W$, $x_i\in V_i$, $x=\sum x_i$, nous dirons que $x_i$ est la \textit{composante} de $x$
dans $V_i$
\end{defin}
% (« dans V_i » : no final period is printed.)

\begin{remarque}
$E$ opère par translation sur l'ensemble $K(E)$ des parties convexes de $E$; on peut définir le
quotient $K(E)/E$; nous dirons qu'une classe $w\in K(E)/E$ est somme de $n$ classes $v_i$ si l'on
peut trouver $W$, $V_i$ de classes $w$, $v_i$ tels que $W$ soit la somme des $V_i$ : dans ce cas,
à tout système de $n$ convexes $V'_i$ de classes $v_i$ correspond un convexe et un seul $W'$ de
classe $w$, qui en est la somme; à tout convexe $W'$ de classe $w$, on peut (et ce d'une infinité
de façons si $n\neq 1$) faire correspondre une famille de $n$ convexes $V'_i$ de classe $v_i$ dont
$W'$ soit la somme.
\end{remarque}

On a le résultat plus précis suivant :

\begin{prop}[1]
Soient $w$, $v_i\in K(E)/E$; $\sum_{i=1}^{n}v_i=w$ et $W$ de classe $w$ contenant l'origine de
$E$; alors il existe un système et un seul de $n$ convexes $V_i$ de classes $v_i$ contenant
chacun l'origine, et tels que
\[ W=\sum_{i=1}^{n}V_i. \]
\end{prop}

\begin{proof}
En effet, supposons déterminés des $V'_i$ tels que
\[ W=\sum_{i=1}^{n}V'_i, \]
On a une décomposition unique de $O$ comme somme de $n$ vecteurs $x'_i\in V_i$. Les seuls convexes
$V_i$, de somme $W$, contenant l'origine et de classe $v_t$ ne sont autres que les
\[ V'_i+\{-x'_i\}. \]
% [sic] « x'_i ∈ V_i » : read « x'_i ∈ V'_i ». [sic] « de classe v_t » : read « v_i ».
\end{proof}

Dans ce qui suit, nous considérerons toujours, pour la commodité des démonstrations, les
décompositions de convexes contenant l'origine en somme de convexes contenant l'origine.

Aux \bzcf{théorèmes 1}{\ref{bzthm:c3s2t1_conjugues_intersection}} et
\bzcf{2}{\ref{bzthm:c3s2t2_algebre_boole_affine}} qui suivent correspondront des énoncés analogues
pour l'espace $K(E)/E$.

% --- p. 264 ---
\begin{remarque}
Pour $I$ infini, on peut encore définir la notion de convexe $W$ somme d'une famillle de convexes
$V_i$ ($i\in I$) : on supposera que $W$ et les $V_i$ contiennent l'origine, et l'on dira que
$W=\sum_{i\in I}V_i$ si tout $x\in W$ s'écrit d'une façon et d'une seule comme une somme finie,
\[ x=\sum_{i\in N\subset I}x_i\qquad(x\neq 0) \]
(où $N$ désigne une partie finie de $I$).

La composante de $x$ dans $V_j$ ($j\notin N$) est $0$.
\end{remarque}
% [sic] « famillle » : read « famille ». The index « i ∈ I » is printed under V_i.

\begin{remarque}
Soient $W=E$, $V_i=E_i$ sous-espace vectoriel de $E$, $W=\sum V_i$ au sens ci-dessus défini :
alors on a $E=\sum E_i$, somme directe des $E_i$, \textit{pour les sous-espaces vectoriels la
notion de somme de convexes coïncide avec celle de somme directe ordinaire}.
\end{remarque}

\begin{prop}[2]
Soit
\[ W=\sum_{i\in I}V_i\qquad\Bigl(x, y\in W;\ x_i.\ y_i\in V_i;\ x=\sum x_i,\ y=\sum y_i\Bigr). \]
on a $x\,\alpha(W)\,y$ (\bzcf{\textit{cf.} déf. 1}{\ref{bzdef:c3s1d1_preordre_alpha}}) si et
seulement si l'on a pour tout $i$ : $x_i\,\alpha(V_i)\,y_i$.
\end{prop}
% [sic] « x_i. y_i ∈ V_i » : read « x_i, y_i ∈ V_i ».
% « cf. déf. 1 » : printed in § 2, where déf. 1 is the sum of convexes; the relation α(W) is
% § 1, déf. 1, which the \ref names (see NOTES.md, question).

\begin{proof}
La démonstration est évidente à partir des définitions.
\end{proof}

Dans ce qui suit, nous étudierons des corps convexes comme polycylindre, c'est-à-dire comme somme
de deux ou plusieurs convexes contenus dans des sous-espaces supplémentaires.

% « DÉFINITION 2, — » : the heading is printed with a comma.
\begin{defin}[2]
Soient $E$ un espace vectoriel, $W$ un convexe de $E$ contenant l'origine; on dit qu'un
sous-espace vectoriel $L$ de $E$ \textit{fibre} $W$ si l'intersection de tout translaté $L'$ de
$L$ avec $W$ est, soit $\emptyset$, soit un translaté de $L\cap W$.
\end{defin}

\begin{defin}[3]
Soit $W$ un convexe de $E$, engendrant $E$, et contenant l'origine; on dit qu'un sous-espace $L$
de $E$ \textit{fibre $W$ cylindriquement}, ou encore que $L$ est \textit{distingué relativement
à $W$}, s'il existe un sous-espace $L'$ de $E$, supplémentaire de $L$, et tel que
\[ W=L\cap W+L'\cap W. \]
\end{defin}

\begin{remarque}
Si $L$ fibre $W$ (resp. fibre cylindriquement $W$), $L$ fibre (resp. fibre cylindriquement) tout
translaté de $W$ contenant l'origine.
\end{remarque}

\begin{prop}[3]
1\textsuperscript{o} Si $W$ ne contient aucune demi-droite, tout sous-espace $L$ qui fibre $W$ le
fibre cylindriquement;
% --- p. 265 ---

2\textsuperscript{o} Si $W$ ne contient aucune droite, à tout $L$ qui fibre cylindriquement $W$
correspond un $L'$ et un seul qui soit supplémentaire de $L$ et tel que
\[ W=L\cap W+L'\cap W. \]
La correspondance entre $L$ et $L'$ est une involution sur l'ensemble des sous-espaces distingués
relativement à $W$.
\end{prop}

Cette proposition légitimera la

\begin{defin}[4]
On dit que deux sous-espaces $L$ et $L'$ sont \textit{conjugués} relativement au convexe $W$ si
\[ W=L\cap W+L'\cap W. \]
\end{defin}

% The proof of Prop. 3 is printed after Déf. 4. The optional argument (editorial, not printed)
% names its target, to which the ledger would otherwise not attach it.
\begin{proof}[Proof of \ref{bzprop:c3s2p3_fibre_cylindriquement}]
Démontrons maintenant la proposition :

1\textsuperscript{o} Désignons par $U'$ l'ensemble des homologues de l'origine dans tous les
translatés de $U=L\cap W$ contenus dans $W$. Puisque $L$ fibre $W$ par hypothèse, il est clair
que $U+U'=W$ et que $U'$ est convexe. Reste à montrer que le support $L'$ de $U'$ est
supplémentaire du support $L$ de $U$ : cela va résulter de ce que tout vecteur de $W$ s'écrit
d'une manière unique comme somme d'un vecteur de $U$ et d'un vecteur de $U'$.

En effet, si un convexe $C$ ne contient aucune demi-droite, il n'existe pas de translation $t$
autre que la translation nulle, qui transforme $C$ en un convexe $t(C)\subset C$.

Soit $w\in W$ : il existe donc un translaté et un seul $U(w)$ de $U$ tel que
$w\in U(w)\subset W$.

2\textsuperscript{o} Soit $w\in W$ : il existe un translaté et un seul $L(w)$ de $L$ tel que
$w\in L(w)$.

$L(w)\cap W$ est un translaté $U(w)$ de $U=L\cap W$; nous allons montrer qu'il existe une
translation unique faisant passer de $U$ à $U(w)$ : il en résultera qu'un convexe $U'$ de support
$L'$ supplémentaire de $L$ et tel que $U'+U=W$, ne peut être que l'ensemble des homologues de
l'origine dans les translations $U(w)$ de $U$.

Or, si un convexe $C$ ne contient aucune droite, il n'existe pas de translation $t$ non nulle qui
transforme $C$ en un convexe $t(C)=C$.

Soit $w\in W$, il existe une translation et une seule qui transforme $U$ en $U(W)$.
% [sic] « U(W) » : read « U(w) ».
\end{proof}

\begin{theoreme}[1]
Soient $W$ un convexe de $E$ contenant l'origine $O$, engendrant $E$, et ne contenant aucune
droite, $L$ et $L'$, $D$ et $D'$ deux couples de sous-espaces distingués conjugués relativement à
$W$. Alors $L\cap D$ et $L'\cap D$ sont conjugués relativement à $W\cap D$.
\end{theoreme}

% Proof of Theorem 1, first part; it is interrupted by the LEMME and its proof, and resumes
% after them. Optional arguments editorial and not printed.
\begin{proof}[Proof of \ref{bzthm:c3s2t1_conjugues_intersection}]
Désignons par $P_L$ (resp. $P_{L'}$, $P_D$, $P_{D'}$) la projection sur $L$ (resp. $L'$, $D$, $D'$)
parallèlement à $L'$ (resp. $L$, $D'$, $D$); la démonstration repose sur le
\end{proof}

% --- p. 266 ---
\begin{lemme}
Soient $a$ et $b$ deux éléments tels que $b$ et $a+b$ appartiennent à $W$. Alors
\[ u=P_L(a)+b\in W. \]
\end{lemme}

\begin{proof}
En effet, on a
\[ \begin{aligned} &P_L(u)=P_L(a)+P_L(b)=P_L(a+b)\in W,\\ &P_{L'}(u)=P_{L'}(b)\in W,
\end{aligned} \]
d'où
\[ u=P_L(u)+P_{L'}(u)\in W. \]
\end{proof}

% Proof of Theorem 1, second part (after the LEMME); NOT a proof of the LEMME.
\begin{proof}[Proof of \ref{bzthm:c3s2t1_conjugues_intersection}]
Soit $a\in W\cap D$, nous allons montrer que $P_L(a)$ et $P_{L'}(a)$ appartiennent à $D$, ce qui
prouvera que
\[ W\cap D=P_L(W\cap D)+P_{L'}(W\cap D) \]
d'où le théorème.

On a
\[ P_{D'}P_L(a)+P_{D'}P_{L'}(a)=P_{D'}(a)=0. \]

Nous allons montrer que $D'\cap W$ est stable par les translations opposées $P_{D'}P_L(a)$ et
$P_{D'}P_{L'}(a)$, d'où résultera la nullité de ces translations.

En effet, soit $w\in D'\cap W$; on a
\[ a+w\in D\cap W+D'\cap W=W, \]
d'où, d'après \bzcf{le lemme}{\ref{bzlem:c3s2l1_projection_somme}},
\[ P_{L'}(a)+w\in W \]
et, en appliquant \bzcf{le lemme}{\ref{bzlem:c3s2l1_projection_somme}} cette fois à $P_{D'}$,
\[ P_{D'}P_{L'}(a)+w\in D'\cap W. \]
On voit de même que
\[ P_{D'}P_L(a)+w\in D'\cap W, \]
ce qui achève la démonstration.
\end{proof}

\begin{corollaire}
Soient $W$ un convexe de $E$ contenant $O$, engendrant $E$, ne contenant aucune droite, $(L, L')$,
$(D, D')$ deux couples de sous-espaces conjugués relativement à $W$, $P_L$ (resp. $P_{L'}$,
$P_D$, $P_{D'}$) la projection sur $L$ (resp. $L'$, $D$, $D'$) parallèlement à $L'$ (resp. $L$,
$D'$, $D$) : $P_L$, $P_{L'}$, $P_D$, $P_{D'}$ commutent deux à deux, et l'on a
\[ \text{Identité}=P_LP_D+P_LP_{D'}+P_{L'}P_D+P_{L'}P_{D'}. \]
\end{corollaire}

\begin{proof}
D'après \bzcf{le théorème 1}{\ref{bzthm:c3s2t1_conjugues_intersection}}, tout point de
$V=D\cap W$ s'écrit d'une façon et d'une seule comme somme d'un point de $L\cap D\cap W$ et d'un
point de $L'\cap D\cap W$. De même, tout point de $V'$ s'écrit d'une manière et d'une seule comme
% --- p. 267 ---
somme d'un point de $L\cap D'\cap W$ et d'un point de $L'\cap D'\cap W$. Tout point de $W$ s'écrit
de manière unique comme somme de points de
\[ L\cap D\cap W;\qquad L'\cap D\cap W;\qquad L\cap D'\cap W;\qquad L'\cap D'\cap W. \]
La projection correspondante sur $L\cap D\cap W$ n'est autre que $P_LP_D$ ou $P_DP_L$ par raison
de symétrie. D'où le corollaire.
\end{proof}
% « V' » as printed (read D' ∩ W; V' is not introduced).

De là on déduit aussitôt le

\begin{theoreme}[2]
Soit $W$ un convexe de $E$ contenant l'origine, engendrant $E$ et ne contenant aucune droite :
l'ensemble des sous-espaces distingués relativement à $W$ peut être considéré comme une algèbre
de Boole, dont l'élément $1$ est $E$ et l'élément $0$ l'origine (considéré comme sous-espace de
dimension $0$); on prend comme complémentaire le conjugué; comme \textup{inf} l'intersection,
comme \textup{sup} le produit.
\end{theoreme}
% (No proof of Theorem 2 is printed beyond « De là on déduit aussitôt le ».)

\begin{corollaire}
Si $E$ est de dimension finie $n$, l'algèbre de Boole des sous-espaces distingués relativement à
$W$ n'est autre que l'algèbre des parties d'un ensemble fini dont le cardinal est compris entre
$1$ et $n$ ($1$ correspond à la sphère, $n$ au cube).

Tout sous-espace distingué est obtenu d'une façon et d'une seule comme somme d'une famille de
sous-espaces distingués simples [c'est-à-dire minimaux parmi ceux différents de $0$].
\end{corollaire}
% (No proof of the corollary is printed.)

Nous allons maintenant étudier la notion de somme pour les convexes de la sphère projective
$S(E)$ associée à l'espace vectoriel $E$.

\begin{defin}[5]
Soient $w$ un convexe de $S(E)$, de noyau vide, $\{v_i\}$ une famille finie de convexes de $S(E)$;
$w$ est dit \textit{somme} des convexes $v_i$ si les $F(v_i)$ sont convexes et si $F(w)$ est
somme des convexes $F(v_i)$.

En particulier, soient $w$, $v_1$, $v_2$ trois convexes non vides de $S(E)$. Soit $l_1$, $l_2$
les supports de $v_1$, $v_2$, c'est-à-dire les plus petites variétés linéaires contenant $v_1$ et
$v_2$ respectivement. $w$ est somme de $v_1$ et $v_2$ si $l_1\cap l_2$ est vide et si $w$ est
réunion de l'ensemble de tous les segments ouverts joignant un point de $v_1$ à un point de
$v_2$. On convient que $w$ est somme de $w$ et du convexe vide.
\end{defin}

\begin{defin}[6]
Soit $w$ un convexe de $S(E)$ dont le support est $S(E)$ et le noyau $\emptyset$. Un sous-espace
$l$ est \textit{distingué relativement à $w$} s'il existe un sous-espace $l'$ tel que $w$ soit
somme de deux convexes $v$ et $v'$ de supports $l$ et $l'$.

La donnée du sous-espace distingué $l$ détermine complètement le sous-espace distingué $l'$; $l$
et $l'$ sont dits coujugués relativement à $w$.
% [sic] « coujugués » : read « conjugués ».
\end{defin}

\begin{theoreme}[3]
Soit $w$ un convexe de $S(E)$ de support $S(E)$, de noyau vide; l'ensemble des sous-variétés
distinguées relativement à $w$ peut être considéré comme une algèbre de Boole dont l'élément
\textup{(}$1$\textup{)} est $S(E)$ et l'élément \textup{(}$0$\textup{)} la partie vide. On prend
comme complémentaire la conjuguée, comme \textup{inf} l'intersection, comme \textup{sup}
l'enveloppe convexe.
\end{theoreme}

% --- p. 268 ---
\begin{corollaire}
Si $E$ est de dimension finie $(n+1)$, l'algèbre de Boole des sous-variétés distinguées
relarivement à $w$ n'est autre que l'algèbre des parties d'un ensemble fini dont le cardinal est
compris entre $1$ et $(n+1)$. [$1$ correspond au cas où il n'existe pas de décomposition non
triviale, $(n+1)$ au cas du $n$-simplexe ouvert de la sphère à $n$ dimensions $S(E)$].
% [sic] « relarivement » : read « relativement ».

Toute sous-variété distinguée est obtenue d'une façon et d'une seule comme enveloppe convexe d'une
famille de sous-variétés distinguées simples, c'est-à-dire minimales parmi celles distinctes de la
partie vide.
\end{corollaire}

% This paragraph justifies Déf. 6 and proves Theorem 3 and its corollary at once; as a `proof`
% without optional argument the ledger attaches it to the corollary (the nearest statement).
\begin{proof}
Pour justifier \bzcf{la définition 6}{\ref{bzdef:c3s2d6_distingue_projectif}} et démontrer
\bzcf{le théorème 3}{\ref{bzthm:c3s2t3_algebre_boole_projective}} et son corollaire, il suffit
de se reporter aux énoncés correspondants formulés dans le cas affin.
\end{proof}

\begin{prop}[4]
Soient $w$ un ouvert convexe de $S(\R^{n+1})=S(n)$, de noyau vide, $l_i$ les $p$ sous-variétés
distinguées simples relatives à $w$, $v_i$ les convexes de $l_i$ dont $w$ est la somme; on peut
écrire avec les notations du chapitre 1, paragraphes 1 et 2
\[ w=F^{-1}(F(v_1)\times\ldots\times F(v_p)) \]
tout automorphisme projectif de $w$ est un automorphisme compatible au sens du \bzcf{chapitre 1,
paragraphe 3, définition 2}{\ref{bzdef:c1s3d2_automorphisme_compatible}}.
\end{prop}

\begin{proof}
En effet, un automorphisme projectif de $w$ est un élément de $GP(n, \R)$ qui permute entre elles
les variétés distinguées simples.
\end{proof}

\begin{prop}[5]
Avec les notations de \bzcf{la proposition 4}{\ref{bzprop:c3s2p4_automorphisme_compatible}}, $w$
est balayable (resp. homogène) si et seulement si les $v_i$ le sont.
\end{prop}

\begin{proof}
Cette proposition résulte de \bzcf{la précédente}{\ref{bzprop:c3s2p4_automorphisme_compatible}}
et de \bzcf{la proposition 4 du chapitre 1, paragraphe 3}{\ref{bzprop:c1s3p4_reciproque}}.
\end{proof}

\begin{corollaire}
Soit $w$ un ouvert convexe de $S(\R^{n+1})$ de noyau vide, somme de $p$ convexes $v_1$, \ldots,
$v_i$, \ldots, $v_p$. $w$ est balayable (resp. homogène) si et seulement si les $v_i$ le sont.
\end{corollaire}

\begin{proof}
\textsc{Démonstration.} --- On se ramène à une décomposition de $w$ en somme de convexes contenus
dans des variétés distinguées simples et l'on applique \bzcf{la proposition
précédente}{\ref{bzprop:c3s2p5_balayable_somme}}.
\end{proof}

% « 3. Dualité des corps convexes — » : no period is printed before the dash.
\subsection*{3. Dualité des corps convexes.}

Nous allons définir cette dualité en introduisant successivement les hypothèses nécessaires :

\textit{a}. Soient $E$, $E'$ deux espaces vectoriels. Supposons définie sur $E\times E'$ une
forme bilinéaire que nous noterons
\[ \langle m, m'\rangle,\qquad m\in E,\qquad m'\in E' \]
Nous allons définir au moyen de cette forme bilinéaire une application $D$ de $K(S(E))$ dans
$K(S(E'))$ et une application $D'$ de $K(S(E'))$ dans $K(S(E))$.
% --- p. 269 ---
Soit $W\in K(S(E))$; nous poserons (pour la définition de $s$, \textit{voir} \bzcf{\S\,1,
chap. 1, déf. 3}{\ref{bzdef:c1s1d3_sphere}})
\[ D(W)=s\{m'\mid\langle m, m'\rangle\geqq 0,\ \forall m\in F(W);\ m'\neq 0\}. \]
et, de même pour $W'\subset K(S(E'))$,
\[ D'(W')=s\{m\mid\langle m, m'\rangle\geqq 0,\ \forall m'\in F(W');\ m\neq 0\}. \]
% [sic] « W' ⊂ K(S(E')) » : read « W' ∈ K(S(E')) ».
Il est clair que $F\circ D(W)$ est un cône convexe de sommet origine dans $E'$ et l'on a bien
défini une application de $K(S(E))$ dans $K(S(E'))$.

Soit $W_1\subset W_2$ deux éléments de $K(S(E))$; on a $D(W_2)\subset D(W_1)$.

\textit{b}. Supposons maintenant que $E$ (resp. $E'$) soit muni d'une topologie localement
convexe séparée et que la forme bilinéaire introduite en (\textit{a}) soit séparément continue en
chacune des variables.

Désignons par $C(S(E))$ l'ensemble des convexes fermés de la sphère $S(E)$ d'un espace vectoriel
topologique $E$ : $D$ envoie $C(S(E))$ dans $C(S(E'))$ et, de même, $D'$ envoie $C(S(E'))$ dans
$C(S(E))$.

\textit{c}. Supposons vérifiées les hypothèses de (\textit{a}) et (\textit{b}) et supposons de
plus que la forme bilinéaire définie sur $E\times E'$ permette d'identifier $E'$ à l'ensemble des
formes linéaires continues sur l'espace vectoriel topologique $E$. Alors $D'\circ D$ induit
l'identité sur $C(S(E))$.

% The claims of a, b, c are left as free text: an environment `proof` here would be attached by
% the ledger to the corollary of § 2 Prop. 5 (see NOTES.md, editorial choices).
En effet, soit $W\subset C(S(E))$; $F(W)\cup\{O\}$ est un cône convexe fermé de $E$, il est donc
l'intersection des demi-espaces fermés qui le contiennent (\textit{cf.} \textsc{Bourbaki},
\textit{Espaces vectoriels topologiques}, chap. II, \S\,3, corollaire 1 de la proposition 4). On
voit facilement que $F(W)$ est l'intersection d'une famille de demi-espaces fermés limités à un
hyperplan contenant l'origine, c'est-à-dire définis par une équation
\[ \langle m, m'\rangle\geqq 0 \]
[si un demi-espace fermé contient $F(W)$, l'intersection de tous les homothétiques de ce
demi-espace par rapport à l'origine dans un rapport positif contient $F(W)$].
% [sic] « soit W ⊂ C(S(E)) » : read « W ∈ C(S(E)) ».

Le résultat annoncé en résulte.

\textit{d}. Nous pouvons maintenant énoncer la

\begin{prop}[1]
Soient $E$, $E'$ deux espaces vectoriels topologiques semi-réflexifs mis en dualité par une
application bilinéaire notée $\langle m, m'\rangle$, $m\in E$, $m'\in E'$. la correspondance $D$
définie en (\textit{a}) est une application biunivoque de $C(S(E))$ sur $C(S(E'))$, et $D'$ est
l'application réciproque.
\end{prop}
% [sic] « m' ∈ E'. la correspondance » : read « La correspondance » (lower-case after a period).

\begin{proof}
Cela résulte de la
\end{proof}
% (« Cela résulte de la » runs into the heading « PROPOSITION 2. »: Prop. 1 is proved by Prop. 2.)

\begin{prop}[2]
1\textsuperscript{o} Sous les hypothèses de (\textit{b}), on a $D(W)=D(\overline{W})$;

2\textsuperscript{o} Sous les hypothèses de (\textit{d}), on a $D'\circ D(W)=\overline{W}$.
\end{prop}

% --- p. 270 ---
\begin{proof}
\textsc{Démonstration.} --- 1\textsuperscript{o} Il suffit de montrer que
$D(W)\subset D(\overline{W})$ : cela résulte de ce que toute forme linéaire positive ou nulle sur
$F(W)$ est positive ou nulle sur sa fermeture;

2\textsuperscript{o} On a $D(W)=D(\overline{W})$ et $D'\circ D=$ identité pour tout convexe de
$C(S(E))$ tel que $\overline{W}$.
\end{proof}
% « tel que W̄ » as printed (« tel que » = « such as »: D'∘D is the identity on the closed convex W̄).

\begin{prop}[3]
Soient $E$ un espace vectoriel topologique, $E'$ l'espace vectoriel des formes linéaires
continues sur $E$, $\langle M, M'\rangle$ la forme bilinéaire sur $E\times E'$ qui à $M'\in E'$,
$M\in E$ associe la valeur sur $M$ de la forme linéaire $M'$. $D'$ met en correspondance
biunivoque les hémisphères fermés de $S(E)$ et les points de $S(E')$.
\end{prop}

\begin{proof}
En effet, soit $m\in S(E')$; soit $M\in F(m)\subset E'$. $D'M$ n'est autre que l'hémisphère fermé
de $E$, fermeture de l'hémisphère ouvert d'équation $M$ (\bzcf{\textit{cf.} chap. 1, \S\,1,
déf. 5}{\ref{bzdef:c1s1d5_hemisphere}}).
% [sic] « l'hémisphère fermé de E » : read « de S(E) ». « D'M » as printed (read D'(m) = D'({sM})).

Réciproquement, soit $H$ un hémisphère ouvert de $S(E)$, soit $M\in E'$ une équation de $H$, on a
\[ \overline{H}=D'\circ sM. \]
\end{proof}

La proposition ci-dessus permet de donner une définition géométrique de $D$ : c'est l'objet de la

\begin{prop}[4]
Supposons vérifiées les hypothèses de (\textit{d}) : soit $W\in C(S(E))$ : $D(W)$ n'est autre que
l'ensemble des points $M'$ de $S(E')$ tels que $W$ soit contenu dans l'hémisphère fermé
$D(\{M'\})$.
\end{prop}
% (No proof of Prop. 4 is printed.)

\begin{prop}[5]
Plaçons-nous sous les hypothèses de (\textit{a}) : Soit $W$ une variété linéaire; alors $D$ est
une variété linéaire.
\end{prop}
% [sic] « alors D est une variété linéaire » : read « D(W) ».

\begin{proof}
La proposition résulte facilement de la définition de $D$.
\end{proof}

Nous supposerons maintenant que $E=\R^{n+1}$; $S(E)=S(n)$. $\R^{n+1}$ est muni d'un produit
scalaire qui permet de l'identifier à son dual. Soit $M=(m_0, \ldots, m_i, \ldots, m_n)$, avec
$m_i\in\R$, $N=(n_0, \ldots, n_i, \ldots, n_n)$.

On pose $\langle M, N\rangle=\sum_{i=0}^{n}m_in_i$.

\begin{prop}[6]
Soit $L$ une sous-variété linéaire de dimension $p$ de $S(n)$ ($-1\leqq p\leqq n$), $D(L)$ est
une sous-variété linéaire de dimension $(n-p-1)$.
\end{prop}

\begin{proof}
En effet, $F(L)\cup\{O\}$ est un sous-espace $\R^{n+1}$ de dimension $(p+1)$, et
$F(DL)\cup\{O\}$ un sous-espace de dimension $(n-p)$.
% « un sous-espace R^{n+1} » as printed (read « un sous-espace de R^{n+1} »).
\end{proof}

\begin{prop}[7]
Soit $W\in C(S(n))$; désignons par $\nu(W)$ la dimension de son noyau, par $\sigma(W)$ la dimension
de son support (c'est-à-dire de la plus
% --- p. 271 ---
petite variété linéaire contenant $W$). $D$ transforme le noyau (resp. support) de $W$ dans le
support (resp. noyau) de $D(W)$ et l'on a
\[ \begin{aligned} &\sigma(W)+\nu(D(W))=n-1,\\ &\nu(W)+\sigma(D(W))=n-1. \end{aligned} \]
\end{prop}

\begin{proof}
la proposition résulte facilement de ce que $D$ inverse les relations d'inclusion et de \bzcf{la
proposition précédente}{\ref{bzprop:c3s3p6_dual_sous_variete}}.
% [sic] « la proposition » : lower-case at the start of a sentence (read « La proposition »).
\end{proof}

\begin{prop}[8]
Le dual d'un corps convexe de $S(n)$ est un convexe.
% [sic?] « est un convexe » : the proof (Prop. 7, σ = n and ν = -1 for the dual) gives « est un
% corps convexe », which is how Prop. 8 is used; see NOTES.md.
\end{prop}

\begin{proof}
Un corps convexe de $S(n)$ (\bzcf{\textit{cf.} chap. 1, \S\,1,
déf. 11}{\ref{bzdef:c1s1d11_corps_convexe}}) est un convexe formé $W$ tel que $\sigma(W)=n$;
$\nu(W)=-1$.
% [sic] « un convexe formé » : read « fermé ».

Les corps convexes ne sont autres que les convexes fermés d'intérieur non vide, et de noyau vide.

La proposition résulte immédiatement de \bzcf{la proposition
7}{\ref{bzprop:c3s3p7_noyau_support}}.
\end{proof}

\begin{prop}[9]
Soit $A\in C(S(n))$ un convexe fermé de $S(n)$, $H=D(\{M\})$ un hémisphère fermé de $S(n)$, les
deux propriétés suivantes sont équivalentes :
\[ \begin{aligned} &1^{\circ}\quad A\subset\mathring{H};\\ &2^{\circ}\quad M\in\intr{DA}.
\end{aligned} \]
\end{prop}

\begin{proof}
La sphère de rayon $1$ de $\R^{n+1}$ est une section de $\R^{n-1}-\{O\}$ au-dessus de $S(n)$.
% [sic] « R^{n-1} - {O} » : read « R^{n+1} - {O} ».

Pour démontrer la proposition, nous identifierons $S(n)$ avec cette sphère munie de sa métrique
géodésique.

La propriété 1\textsuperscript{o} peut s'énoncer ainsi : $A$ est contenu dans une boule $B$ de
centre $M$ et de rayon $\dfrac{\pi}{2}-\varepsilon$, $\varepsilon>0$. La propriété
1\textsuperscript{o} est équivalente à
\[ D(B)\subset A. \]
% [sic] « D(B) ⊂ A » : read « D(B) ⊂ D(A) » (D reverses inclusions; D(B) is a ball about M).
Or $(D(B)$ n'est autre que la boule du centre $M$ et de rayon $\varepsilon$. Donc, la propriété
1\textsuperscript{o} est équivalente à la propriété 2\textsuperscript{o}.
% [sic] « Or (D(B) » : stray parenthesis. [sic] « la boule du centre M » : read « de centre M ».
\end{proof}

\begin{corollaire}
Si $A$ est un corps convexe de $S(n)$, on peut trouver un hémisphère ouvert de $S(n)$ contenant
$A$.
\end{corollaire}
% (No proof of the corollary is printed.)

\begin{prop}[10]
Soient $A$, $B$ des convexes fermés : $A$, $B\in C(S(n))$.
\[ A\subset\mathring{B}\quad\text{entraîne}\quad D(B)\subset\intr{D(A)} \]
et réciproquement.
\end{prop}

\begin{proof}
Puisque $D\circ D=$ Identité, la réciproque se change en la proposition en posant
\[ D(A)=B',\qquad D(B)=A'. \]

% --- p. 272 ---
Démontrons que $A\subset\mathring{B}$ entraîne $D(B)\subset\intr{D(A)}$. On a
\[ \begin{aligned} D(A)&=\bigcap_{M\in A}D(\{M\}),\\
\intr{D(A)}&=\bigcap_{M\in A}\intr{D(\{M\})}. \end{aligned} \]
Pour démontrer la proposition, il suffit donc d'établir que pour tout $M\in A$,
$D(B)\subset D(\{M\})$. Or cela résulte de \bzcf{la proposition
précédente}{\ref{bzprop:c3s3p9_hemisphere_dual}}.
% [sic?] « D(B) ⊂ D({M}) » : printed without the ring of interior; the display just above needs
% « D(B) ⊂ interior of D({M}) ».
\end{proof}

Quand, au chapitre suivant (\textit{cf.} chap. 4. \S\,1, prop. 9 et 11), nous aurons défini la
topologie de $E(n)$, de \bzcf{la proposition 10}{\ref{bzprop:c3s3p10_interieur_dual}} résultera
la
% « chap. 4. § 1 » as printed (read « chap. 4, § 1 »). Chapter 4 is not transcribed: the
% reference is not marked.

\begin{prop}[11]
$D$ est un automorphisme topologique de $E(n)$.
\end{prop}
% (No proof of Prop. 11 is printed beyond the sentence above.)

\begin{remarque}
Puisque $\R^n$ est canoniquement identifié à son dual, la représentation duale de $GL(n, \R)$ est
un automorphisme de $GL(n, \R)$ : la transposition. Cet automorphisme induit un automorphisme de
$GS(n-1, \R)$, que nous noterons $\Delta$; $\Delta$ peut s'exprimer au moyen de la dualité $D$
des parties convexes fermées de $S(n-1)$.

Soit $g\in(GS(n-1, \R)$, $\Delta(g)\in GS(n-1, \R)$; $W\in C(S(n-1))$; on a
\[ \Delta(g)\,W=D\circ g\circ DW. \]
% [sic] « g ∈ (GS(n-1, R) » : stray parenthesis.
\end{remarque}

\begin{proof}
En effet, cette formule est équivalente à la définition de $\Delta$, quand $W$ est un point ou un
hémisphère.
\end{proof}

\subsection*{4. Contingent des ensembles convexes.}

\begin{defin}[1]
Soient $E$ un espace vectoriel, $W\in K(S(E))$, $a\in W$, $a'$ le point diamétralement opposé à
$a$. On appelle \textit{cône d'appui} en $a$ à $W$ le plus petit convexe de $S(E)$ contenant $W$
et $a'$ : on le note $W_a$.
\end{defin}

On peut donner de cette définition deux autres énoncés équivalents.

\begin{defin}[1$'$]
$W_a$ est la réunion des demi-droites de $S(E)$ qui joignent $a$ au point diamétralement opposé
$a'$ et rencontrent $W$ en un point distinct des extrémités $a$ et $a'$.
\end{defin}

\begin{defin}[1$''$]
Soit $A$ un point de $F(W)$ se projetant sur $W$ suivant $a$ : $F(W_a)\cup\{O\}$ est l'ensemble
des vecteurs de $E$ qui peuvent s'écrire comme somme d'un vecteur de $F(W)$ et d'un vecteur
multiple de $A$ par un nombre réel
\[ F(W_a)\cup\{O\}=\{M : \exists\lambda\in\R;\ \exists P\in F(W);\ M=\lambda A+P\}. \]
\end{defin}

\begin{defin}[2]
On appelle contingent à un convexe fermé $W$ en un de ses points $a$ la fermeture du
cône d'appui $W_a$.
\end{defin}

% --- p. 273 ---
\begin{prop}[1]
Soient $W\in K(S(E))$, $a$, $b\in W$, $a'$, $b'$ les points diamétralement opposés à $a$ et $b$.

Si $b\neq a$, $b\neq a'$, les quatre propriétés suivantes sont équivalentes :

(1) $a\,\alpha(W)\,b$;

(2) Sur la droite qui joint $a$ et $b$, chacune des deux demi-droites délimitées par $a$ et $a'$
rencontre $W$;

(3) $b'\in W_a$;

(4) $W_b\subset W_a$.

Si $b=a$ ou $b=a'$, les propriétés (1), (3), (4) sont équivalentes.
\end{prop}

\begin{proof}
Supposons d'abord $b\neq a$, $b\neq a'$. L'équivalence des propriétés (1) et (2) résulte
immédiatement de la définition de $\alpha(W)$ (\bzcf{\textit{cf.} \S\,2
déf. 3}{\ref{bzdef:c3s2d3_fibre_cylindriquement}}).
% [sic] « cf. § 2 déf. 3 » : α(W) is § 1, déf. 1 (and α(k) on the sphere § 1, déf. 4); § 2,
% déf. 3 is « L fibre W cylindriquement ». Following the convention, the \ref names the PRINTED
% § 2, déf. 3; see NOTES.md (to be ruled).

L'équivalence des propriétés (2) et (3) apparaît si l'on utilise \bzcf{la définition
(1$'$)}{\ref{bzdef:c3s4d1prime_cone_appui_demi_droites}} de $W_a$.

L'équivalence des propriétés (3) et (4) apparaît si l'on utilise \bzcf{la définition
(1)}{\ref{bzdef:c3s4d1_cone_appui}} de $W_a$.

Dans les cas $b=a$ ou $b=a'$, l'assertion de la proposition est évidente.
\end{proof}

\begin{prop}[2]
Soient $E$, $E'$ deux espaces vectoriels. On suppose définie une forme bilinéaire sur le produit
$E\times E'$ [\textit{cf.} \S\,3, \textit{a})]. Soit $W\subset K(S(E))$, $a\in W$, on a
\[ D(W_a)\subset D(W), \]
et si
\[ m'_1\in D(W_a),\qquad m'_2\in D(W)\qquad\text{et}\qquad m'_1\,\alpha(D(W))\,m'_2, \]
alors
\[ m'_2\in D(W_a). \]
\end{prop}
% [sic] « W ⊂ K(S(E)) » : read « W ∈ K(S(E)) ». « cf. § 3, a) » names the hypothesis a of § 3,
% which is free text: not marked.

\begin{proof}
\textsc{Démonstration.} --- Soit $A$ un point de $F(W)$ se projetant en $a$. On peut trouver
$M'_1$, $M'_2$ se projetant en $m'_1$, $m'_2$ tels que $M'_1=M'_2+M'_3$, où
$M'_3\in F\circ D(W)$. $F\circ D(W_a)$ peut être ainsi défini
\[ F\circ D(W_a)\cup\{O\}=\{M'\mid\langle AM'\rangle=0;\ \langle MM'\rangle\geqq 0,\ \forall M\in
F(W)\}. \]
Rappelons la définition de $F\circ D(W)$
\[ F\circ D(W)\cup\{O\}=\{M'\mid\langle MM'\rangle\geqq 0,\ \forall M\in F(W)\}. \]
Pour démontrer la proposition, il suffit d'établir que $\langle AM''_2\rangle=0$. Or on a
\[ \begin{gathered} \langle AM'_2\rangle\geqq 0,\\ \langle AM'_3\rangle\geqq 0,\\
\langle AM'_2\rangle+\langle AM'_3\rangle=\langle AM'_1\rangle=0. \end{gathered} \]
% [sic] « ⟨A M''_2⟩ » : read « ⟨A M'_2⟩ ».
Il en résulte
\[ \langle AM'_2\rangle=0, \]
d'où la proposition.
\end{proof}

% --- p. 274 ---
Nous supposerons maintenant que l'espace $E$ est $\R^{n+1}$, canoniquement identifié à son dual au
moyen de la forme $\sum_i x_iy_i$.

\begin{defin}[3]
Soient $m$, $p$ deux points de $S(n)$; $M$, $P$ deux points de $\R^{n+1}-\{O\}$ se projetant sur
$m$ et $p$. Si $\langle M, P\rangle=0$, nous écrirons encore $\langle m, p\rangle=0$. Il est clair
que cette définition ne dépend pas du choix de $M$ et de $P$ et qu'elle est symétrique en $m$ et
$p$.

On peut encore l'énoncer :

$\langle m, p\rangle=0$ si et seulement si $D(\{p\})$ contient $m$ et le point diamétralement
opposé à $m$.
\end{defin}

\begin{prop}[3]
Soit $a$ un point de $S(n)$; $W\in K(S(n))$; alors
\[ D(W_a)=\{m\mid m\in D(W),\ \langle m, a\rangle=0\}. \]
\end{prop}

\begin{proof}
En effet,
\[ D(W_a)=D(\overline{W_a}) \]
et le contingent $\overline{W_a}$ n'est autre que l'intersection de la famille des hémisphères
fermés contenant $W$ et le point diamétralement opposé à $a$.
\end{proof}

\begin{prop}[4]
Soient $m$, $m'$, $p$, $p'$ quatre points de $S(n)$, entre lesquels on suppose vérifiées les
relations
\[ m\,\alpha(W)\,m';\quad p\,\alpha(D(W))\,p';\quad\langle m, p\rangle=0. \]
Alors
\[ \langle m', p'\rangle=0. \]
\end{prop}

\begin{proof}
Par raison de symétrie, il suffit de démontrer
\[ \langle m', p\rangle=0. \]
Or cela résulte immédiatement des \bzcf{propositions 1}{\ref{bzprop:c3s4p1_alpha_cone_appui}} et
\bzcf{3}{\ref{bzprop:c3s4p3_dual_cone_appui}}.
\end{proof}

\bzcf{La proposition 3}{\ref{bzprop:c3s4p3_dual_cone_appui}} permet d'ailleurs d'établir
l'équivalence entre \bzcf{la proposition 2}{\ref{bzprop:c3s4p2_dual_cone_appui_sature}} et
l'assertion $\langle m, p'\rangle=0$.

Sur $S(n)$, \bzcf{la proposition 2}{\ref{bzprop:c3s4p2_dual_cone_appui_sature}} peut donc être
démontrée comme un corollaire de \bzcf{la proposition 1}{\ref{bzprop:c3s4p1_alpha_cone_appui}}.

\begin{defin}[4]
Soient $F$ une facette de $W$, $F'$ une facette de $D(W)$, nous dirons qu'est vérifiée la
relation
\[ \langle F, F'\rangle=0 \]
si pour tout point $m$ de $F$ et tout point $m'$ de $F'$ on a
\[ \langle m, m'\rangle=0. \]
\end{defin}

D'après \bzcf{la proposition précédente}{\ref{bzprop:c3s4p4_orthogonalite_facettes}}, s'il existe
$m$ dans $F$ et $m'$ dans $F'$ tels que $\langle m, m'\rangle=0$, on a $\langle F, F'\rangle=0$.

% --- p. 275 ---
\begin{prop}[5]
Soient $F_1$, $F_2$ (resp. $F'_1$, $F'_2$) deux facettes de $W$ [resp. de $D(W)$] telles que soient
vérifiées les relations
\[ \begin{gathered} F_1\,\alpha(W)\,F_2;\quad F'_1\,\alpha(D(W))\,F'_2;\\ \langle F_1,
F'_1\rangle=0. \end{gathered} \]
Alors, on a
\[ \langle F_2, F'_2\rangle=0. \]
\end{prop}

\begin{proof}
La proposition 5 est un corollaire de \bzcf{la proposition
4}{\ref{bzprop:c3s4p4_orthogonalite_facettes}}.
\end{proof}

\begin{prop}[6]
Soient $W\in K(S(n))$, $F$ une facette de $W$, $a\in F$, $D(W_a)$ est la fermeture d'une facette
de $D(W)$ qui ne dépend pas du choix de $a$ et que nous noterons $\Dc(F)$.
\end{prop}

\begin{proof}
En effet, considérons un point $a'$ intérieur au convexe $D(W_a)$ muni de la topologie induite par
celle de son support, $\Dc(F)$ n'est autre que la facette de $a'$.
\end{proof}

\begin{prop}[7]
Soient $F$ une facette de $W$, $F'$ une facette de $D(W)$ telles que
\[ \langle F, F'\rangle=0. \]
Alors, on a
\[ \Dc(F)\,\alpha(D(W))\,F'. \]
\end{prop}
% (No proof of Prop. 7 is printed.)

Nous dirons que $\Dc(F)$ est la facette minimale (au sens de $\alpha$) parmi celles $F'$ qui
vérifient $\langle F, F'\rangle=0$. Du point de vue ensembliste, notons que $\overline{\Dc(F)}$
contient toutes les facettes $F'$ telles que $\langle F, F'\rangle=0$.

La correspondance $\Dc$ n'est pas nécessairement biunivoque, mais voici quelques propriétés de
$\Dc$ :

% « PROPOSITION 8. » : the heading is printed without the dash, the statement displayed.
\begin{prop}[8]
\[ \Dc\circ\Dc\circ\Dc=\Dc. \]
\end{prop}

\begin{proof}
\textsc{Démonstration.} --- $\Dc\circ\Dc(F)$ n'est autre que la facette de $W$ minimale au sens de
$\alpha$ parmi celles $F_i$ qui vérifient $\langle F_i, \Dc(F)\rangle=0$.

$\Dc\circ\Dc\circ\Dc(F)$ est la facette de $D(W)$ minimale parmi celles $F'$ qui vérifient
$\langle\Dc\circ\Dc(F), F'\rangle=0$.

Or, $\Dc(F)$ est minimale parmi celles $F'$ vérifiant $\langle F, F'\rangle=0$, donc \textit{a
fortiori} parmi celles vérifiant $\langle\Dc\circ\Dc(F), F'\rangle=0$.
\end{proof}

\begin{prop}[9]
Une facette $F$ vérifie $\Dc\circ\Dc(F)=F$ si et seulement si $\overline{F}$ est l'intersection de
$W$ avec un hyperplan.
\end{prop}

\begin{proof}
\textsc{Démonstration.} --- Si $\Dc\circ\Dc(F)=F$, soit $p\in\Dc(F)$, $\overline{F}$ n'est autre
que l'intersection de $W$ avec l'hyperplan frontière de $D(\{p\})$.

Réciproquement, si $\overline{F}=W\cap h$, où $h$ est l'hyperplan frontière de $D(\{p\})$,
$p\in W$, $F$ est transformée par $\Dc$ de la facette $F'$ de $p$, donc
\[ \Dc\circ\Dc(F)=\Dc\circ\Dc\circ\Dc(F')=\Dc(F')=F. \]
% [sic?] « p ∈ W » : p, the pole of a hyperplan meeting W in F̄, lies in D(W) (read « p ∈ D(W) »).
\end{proof}

% --- p. 276 ---
Donnons un exemple de corps convexe possédant une facette qui n'est pas l'intersection du corps
avec un hyperplan.

Nous nous placerons sur $S(2)$ identifié à la sphère unité de $\R^3$ et muni de la métrique
géodésique.

Soient $a$, $b$ deux points dont la distance est $1$; nous prendrons pour $W$ le plus petit
convexe fermé contenant les deux cercles de centres $a$ et $b$ et de rayon $\dfrac{1}{4}$.

La frontière de $W$ se compose de deux arcs de cercle et de deux segments de droite, soit $p$ un
des quatre points où se raccordent un segment et un arc, $\{p\}$ une facette qui ne peut être
obtenue comme intersection de $W$ et d'une droite.
% (The example is printed without an EXEMPLE heading: free text, no environment.)

\begin{defin}[5]
Une facette $F$ d'un convexe fermé $W$ est dite \textit{conique} si le noyau du contingent en un
point de cette facette n'est autre que le support de la facette.
\end{defin}

\begin{prop}[10]
Soient $W$ un corps convexe, $F$ une facette conique de $W$, $L$ son support; le contingent à $W$
en un point de $F$ s'identifie à un corps convexe de $S(n)/L$.
\end{prop}

\begin{proof}
Parce que $W$ est d'intérieur non vide, il en est de même du contingent considéré comme partie de
$S(n)/L$; que le noyau de cette partie soit vide résulte de ce que $F$ est conique.
\end{proof}

\begin{defin}[6]
Soit $F$ une facette conique de support $L$ d'un corps convexe de $W$; soit $L'$ une sous-variété
linéaire de $S(n)$ telle que $S(n)$ soit somme de $L$ et de $L'$ au sens de \bzcf{la définition 6
du paragraphe 2}{\ref{bzdef:c3s2d6_distingue_projectif}} [$L\cap L'=E$, $S(n)$ enveloppe convexe
de $L$ et $L'$]. L'intersection de $L'$ avec le contingent à $W$ en un point de $F$ est appelée
base dans $L'$ de la facette conique $F$.
% [sic] « d'un corps convexe de W » : read « d'un corps convexe W ». [sic] « L ∩ L' = E » : read
% « L ∩ L' = ∅ » (S(n) is the sum of L and L', whose supports are disjoint).
\end{defin}

\begin{remarque}
Supposons données $p$ sous-variétés linéaires $L_1$, \ldots, $L_i$, \ldots, $L_p$ dont $S(n)$
soit la somme. Soient $W_1$, \ldots, $W_i$, \ldots, $W_p$ des corps convexes de $L_1$, \ldots,
$L_i$, \ldots, $L_p$ ($W_i$ convexe fermé d'intérieur non vide pour la topologie de $L_i$). La
fermeture de la somme des $W_i$ est un corps convexe de $S(n)$ dont les $W_i$ sont des facettes
coniques.
\end{remarque}

\subsection*{5. Sur les convexes, réunions de deux convexes disjoints.}

Le but de ce paragraphe et de démontrer les \bzcf{théorèmes
1}{\ref{bzthm:c3s5t1_reunion_convexes_affine}} et
\bzcf{1$'$}{\ref{bzthm:c3s5t1prime_reunion_convexes_sphere}} qui concernent l'espace affin et la
sphère projection de dimension finie.
% [sic] « Le but de ce paragraphe et de démontrer » : read « est de démontrer ». [sic] « la sphère
% projection » : read « projective ».

Plusieurs des résultats intermédiaires sont indépendants de la dimension; nous les démontrons donc
sans faire d'hypothèses restrictives.

\begin{defin}[1]
Soit $W$ un convexe dans un espace vectoriel; $V\subset W$ : on dit que $V$ possède relativement à
$W$ la propriété $I$ si $V$ et $W-V$ sont tous deux convexes, et l'on écrira :
$I(V, W)$.
\end{defin}

% --- p. 277 ---
\begin{prop}[1]
Si la réunion $W$ de deux convexes $U$ et $V$ est convexe, l'une des variétés support est contenue
dans l'autre.

D'une façon plus précise : ou bien tout point de $U$ est aligné avec deux points distincts de
$V$, ou bien tout point de $V$ est aligné avec deux points distincts de $U$.
\end{prop}

\begin{proof}
\textsc{Démonstration.} --- Si l'on écarte la première éventualité, il existe $a\in U$ tel que pour
tout point $b\in V$ le segment demi-ouvert $)a, b]$ soit contenu dans $U$; donc $b$ est aligné
avec deux points distincts de $U$.
\end{proof}

\begin{prop}[2]
Soient $W$ un convexe dans un espace vectoriel, $V$ une partie convexe non vide de $W$ telle que
$I(V, W)$; la variété linéaire $L$ engendrée par $V$ est une variété d'appui de $W$.
\end{prop}

\begin{proof}
Soit $B$ un point de $L\cap W$; nous allons montrer qu'il existe un point $A$ de $V$ tel que
$A\,\alpha\,B$ : en effet, $B$ appartient à une sous-variété linéaire de dimension finie engendrée
par des points $A_1$, \ldots, $A_i$, \ldots, $A_n$ de $V$ : si $A$ est un point intérieur à
l'enveloppe convexe des $A_i$ pour la topologie de son support, il est clair que $A\in V$ et
$A\,\alpha\,B$.

Soit maintenant $B\in W$ tel qu'il existe $A\in L\cap W$, avec $A\,\alpha\,B$ : d'après ce qui
vient d'être démontré, on peut supposer que $A\in V$ : pour démontrer la proposition, il faut
établir que $B\in L$; ceci est évident si $B\in V$; nous considérons sur la droite $AB$ un point
$A'$ tel que $A\in)A', B($; on a certainement $A'\in V$ : en effet, si $A'$ appartient à
$W-V$, $(A', B)$ est contenu dans $(W-V)$, puisque cet ensemble est convexe, et $A$ ne peut
appartenir à $V$. Il en résulte que la droite $AB$ est contenue dans $L$.
\end{proof}

\begin{prop}[3]
Soit $L$ une variété d'appui de $W$ : $\complement L\cap W$ est convexe. Soit $C$ une partie
convexe de $L\cap W$ : $\bigl(W\cap\complement L\bigr)\cup C$ est convexe.
\end{prop}

\begin{proof}
Il suffit de montrer que quels que soient
\[ A\in\bigl(W\cap\complement L\bigr)\cup C\qquad\text{et}\qquad B\in W\cap\complement L, \]
l'ensemble $)A, B(\cap L$ est vide : en effet, si $C\in)A, B(\cap L$, on a $C\,\alpha\,B$,
donc, d'après \bzcf{la définition de la variété d'appui}{\ref{bzdef:c3s1d3_variete_appui}} $L$,
$B\in L$.
% « C » names both the convex part of L ∩ W and a point of ]A, B[ ∩ L, as printed.
\end{proof}

\begin{prop}[4]
Soit $V$ une partie convexe non vide de $W$ convexe : il est équivalent de dire que $V$ possède
la propriété $I$ relativement à $W$, ou de dire que le support $L$ de $V$ est une variété d'appui
de $W$ et que $V$ possède la propriété $I$ relativement à $L\cap W$.
\end{prop}
% (No proof of Prop. 4 is printed.)

\begin{prop}[5]
Soient $V$, $W$ convexes de même support $L$, $V$ fermé, $V\subset W$, $I(V, W)$ : alors, ou bien
$V=W$, ou bien le support de $W-V$ est $L$.
\end{prop}

\begin{proof}
Montrons que si $W-V$ contient un point $C$, son support contient tout
% --- p. 278 ---
point $A$ de $V$. En effet, la droite $CA$ coupe $V$ suivant un segment fermé $(A'A'')$ ou une
demi-droite $(A', \infty)$ ne contenant pas $C$ et le plus petit des segments ouverts non vides
$)C, A'($, $)C, A''($ appartient à $W-V$.
\end{proof}

\begin{defin}[2]
Soient $U$, $V$ deux convexes tels que $U\cap V=\emptyset$. On appelle \textit{ensemble de
contact} de $U$ et $V$ l'ensemble $F$ des points $A$ tels qu'il existe deux points $B$ et $C$
alignés avec $A$, et tels que $)B, A(\subset U$; $)A, C(\subset V$ (comme dans \bzcf{\S\,1,
déf. 1}{\ref{bzdef:c3s1d1_preordre_alpha}}, $)A, A($ désigne $\{A\}$).
\end{defin}

\begin{prop}[6]
L'ensemble de contact $F$ de deux convexes est conxexe.
% [sic] « conxexe » : read « convexe ».
\end{prop}

\begin{proof}
Il suffit, pour le démontrer, de considérer deux points $A$ et $A'$ de $F$, et des points $B$, $C$
ou $B'$, $C'$, correspondant à \bzcf{la définition 2}{\ref{bzdef:c3s5d2_ensemble_contact}} :
$ABC$ (resp. $A'B'C'$) alignés; $)B, A(\subset U$; $)A, C(\subset V$; $)B', A'(\subset U$;
$)A', C'(\subset V$. L'enveloppe convexe des segments $)BA($, $)B'A'($ (resp. $)AC($,
$)A'C'($) est contenu dans $U$ (resp. $V$); et le segment $(AA')$ est contenu dans l'ensemble de
contact de ces deux enveloppes convexes, donc dans $F$.
\end{proof}

\begin{prop}[7]
Soient $V$, $W$ deux convexes non vides de même support $L$, qu'on suppose muni de la topologie
localement convexe la plus fine; supposons $V$ fermé, $V\neq W$, $I(V, W)$. Il existe alors un
hyperplan et un seul séparant $V$ de $W-V=U$.
\end{prop}

\begin{proof}
L'ensemble de contact de $V$ et de $(W-V)$ est l'ensemble $F$ des points $A$ de $V$ tels qu'il
existe un $B\in W$ avec $)A, B(\subset\complement V$; si un hyperplan $H$ sépare $V$ de $W-V$, on
doit avoir $F\subset H$. Nous allons montrer que le support $H$ de $F$ est un hyperplan qui sépare
$V$ de $W-V$, et que $H\cap W=F$. Établissons d'abord le théorème dans le cas où $L$ est de
dimension finie : d'après \bzcf{la proposition 5}{\ref{bzprop:c3s5p5_support_complementaire}},
$L$ est engendré par $U$ (resp. $V$) et comme la dimension est finie, $\mathring{U}$ (resp.
$\mathring{V}$) n'est pas vide; il en est de même de l'intérieur de $F$ pour la topologie de son
support $H$.
% « le théorème » as printed: Prop. 7.

Soit $B\in\mathring{U}$; $C\in\mathring{V}$ : $)B, C($ coupe $F$ au point $A$; et il en est de
même de toute parallèle à $BC$ suffisamment voisine de cette droite : la codimension de $H$ est
donc au plus $1$. Or $H$ ne peut être tout l'espace $L$, car l'intérieur de $F$ dans $L$ est
manifestement vide : donc $H$ est un hyperplan.

$H$ qui contient l'ensemble de contact de $U$ et $V$, dont la réunion est convexe, sépare $U$ et
$V$ : car tout segment joignant un point de $U$ à un point de $V$ doit couper $F$, donc $H$.
Enfin, $H\cap W=F$ : soit $A'\in H\cap W$; nous allons montrer que $A'\in F$. En effet, soient
$C\in\mathring{V}$; $B\in\mathring{V}$ : $)A', C(\subset\mathring{V}$, et donc, puisque $V$ est
fermé, $A'\in V$; $)A', B(\subset\mathring{U}$, donc $A'\in F$.
% [sic] « B ∈ V̊ » : read « B ∈ Ů » (then ]A', B[ ⊂ Ů).

Supposons maintenant que $L$ soit de dimension infinie : la démonstration est analogue à celle
du cas fini : il faut essentiellement établir que $H$ est un hyperplan. Montrons que la
co-dimension de $H$ est au plus $1$. Soit $\Delta$ une droite non parallèle à $H$ : il faut
établir que $\Delta$ coupe $H$. Soit $B\in U$, nous nous plaçons dans un sous-espace de dimension
finie, $P$, engendré par $n$
% --- p. 279 ---
points $C_1$, \ldots, $C_n\in V$ et contenant $\Delta$ et $B$ : un tel $P$ existe, parce que $V$
engendre $L$. Dans $P$ on retrouve la situation de dimension finie :
\[ \begin{gathered} P\cap V\quad\text{convexe fermé d'intérieur non vide,}\\
P\cap W=(P\cap V)\cup(P\cap U),\\
P\cap U,\ \text{convexe non vide, et d'intérieur non vide,}\\
\text{d'après \bzcf{la proposition 5}{\ref{bzprop:c3s5p5_support_complementaire}}.}
\end{gathered} \]
Il est clair que l'ensemble de contact $F_P$ de $P\cap V$ et $P\cap U$ dans $P$ est contenu dans
$F$, ensemble de contact de $U$ et $V$ : l'hyperplan $P'$ support de $F_P$ est donc contenu dans
$H$ et il coupe $\Delta$, ne pouvant lui être parallèle.

Pour établir que $H$ est effectivement de co-dimension $1$, il suffit de montrer que $U\cap H$ est
vide; il en résultera de plus que $H\cap W\subset V$. En effet, soit $B\in U$, supposons que
$B\in H$ : nous allons montrer que cette hypothèse a des conséquences absurdes : $B$
appartiendrait à un sous-espace de dimension finie de $H$ engendré par $m$ points $A_1$, \ldots,
$A_m$ de $F$. Ces points $A_i$ sont extrémités de segments $)A_iB_i(\subset V$ : considérons un
sous-espace $P$ de dimension finie contenant $PA_i$ et $PB_i$, donc aussi $B$, et engendré par un
nombre fini de points $C_1$, \ldots, $C_n$ de $V$ : dans $P$ on retrouve encore la situation de
dimension finie : $P$ engendré par $P\cap V$; $P\cap U$ non vide; l'hyperplan de $P$ qui sépare
$P\cap U$ et $P\cap V$ ne peut contenir $B\in P\cap U$, et d'autre part il contient les $A_i$ qui
appartiennent à l'ensemble de contact de $P\cap U$ et $P\cap V$ : $C$ n'appartient donc pas à la
variété engendrée par $PA_i$, contrairement à l'hypothèse.
% « contenant PA_i et PB_i », « la variété engendrée par PA_i » as printed (read: the points
% A_i, B_i). [sic?] « C n'appartient donc pas » : read « B » (no point C is defined here; B is the
% point supposed to lie in H).

Enfin, il est évident que $H\cap W$ n'est autre que $F$ : soit $B\in U$; $A\in H\cap W$, on a
$)A, B(\subset U$ et $A\in V$. Cela achève la démonstration.
\end{proof}

\begin{prop}[8]
Soit $W$ un convexe compact d'un espace vectoriel topologique $E$ localement convexe séparé, $V$
une partie compacte convexe non vide de $W$; alors les deux propriétés ci-dessous sont
équivalentes :

$\mathrm{I}(V, W)$ : $(W-V)$ est convexe;

$\mathrm{II}(V, W)$ : les demi-espaces ouverts contenant $V$ coupent $W$ suivant un système
fondamental de voisinages de $V$ considéré comme partie de $W$.

$\mathrm{II}(V, W)$ peut encore s'énoncer : Pour tout voisinage $O$ de $V$ considéré comme partie
de $E$, il existe un demi-espace ouvert $H$ tel que
\[ \begin{aligned} &1^{\circ}\quad V\subset H;\\ &2^{\circ}\quad H\cap W\subset O\cap W.
\end{aligned} \]
\end{prop}
% « I(V, W) » is printed in roman here, the property I of Déf. 1.

\begin{proof}
Montrons que II implique I. En effet, soient $b$, $b'$ deux points de $W-V$, $O$ un voisinage de
$V$ ne contenant ni $b$ ni $b'$ (par exemple tout l'espace privé des deux points $b$ et $b'$).
Soit $H$ un demi-espace ouvert vérifiant les conditions 1\textsuperscript{o} et
2\textsuperscript{o}; alors, d'après 2\textsuperscript{o},
\[ b\in\complement H,\qquad b'\in\complement H,\qquad\text{d'où}\qquad(b, b')\in\complement H, \]
% --- p. 280 ---
puisque $\complement H$ est convexe; et, de plus, $(b, b')\in W-V$ : en effet, $(b, b')\subset W$
parce que $W$ est convexe, et aucun point de $(b, b')$ n'est dans $V$, sinon il serait dans $H$
d'après 1\textsuperscript{o}.
% « (b, b') ∈ ∁H », « (b, b') ∈ W - V » as printed (read ⊂: (b, b') is a segment).

Montrons que I implique II. Il suffit d'établir que les $H\cap W$ (où $H$ désigne un demi-espace
ouvert contenant $V$) forment une base de filtre : cette base étant formée de voisinages de $V$
relativement à $W$ compact, dont l'intersection est $V$, le filtre sera le filtre de voisinages
de $V$. Soit donc, en passant aux complémentaires, $H_1$ et $H_2$ deux demi-espaces fermés tels
que
\[ H_1\cap V=H_2\cap V=\emptyset; \]
nous allons montrer qu'il existe un demi-espace fermé $H_3$ tel que
\[ H_3\cap V=\emptyset\qquad\text{et}\qquad H_3\cap W\supset(H_1\cap W)\cup(H_2\cap W). \]
Or, $H_1\cap W$ et $H_2\cap W$ sont deux compacts convexes contenus dans $W-V$; leur enveloppe
convexe $C$ est compacte et contenue dans $W-V$ d'après l'hypothèse I; on peut alors séparer
strictement $C$ et $V$ par un hyperplan (\textsc{Bourbaki}, \textit{Espaces vectoriels
topologiques}, chap. II, \S\,3, prop. 4).
\end{proof}

\begin{corollaire}[1]
Soient $W$, $V$ convexes compacts $V\subset W$; si le support $L$ de $V$ est une variété d'appui
de $W$, il est équivalent d'affirmer $\mathrm{II}(V, W)$ ou $\mathrm{II}(V, W\cap L)$
(\bzcf{\textit{cf.} prop. 4}{\ref{bzprop:c3s5p4_propriete_I_appui}}).
\end{corollaire}
% (No proof of Corollaries 1 and 2 is printed.)

\begin{corollaire}[2]
Soit $W$ convexe compact, $L$ une variété d'appui fermée de $W$, alors on a
$\mathrm{II}(L\cap W, W)$.
\end{corollaire}

Des \bzcf{propositions 7}{\ref{bzprop:c3s5p7_hyperplan_separant}},\bzcf{8}{\ref{bzprop:c3s5p8_proprietes_I_II}}
et \bzcf{4}{\ref{bzprop:c3s5p4_propriete_I_appui}} résulte le

\begin{theoreme}[1]
Soit $W$ un convexe compact dans un espace vectoriel de dimension finie, $V$ une partie non vide
compacte convexe de $W$; alors les trois propriétés ci-dessous sont équivalentes :

\textup{I}\,$(V, W)$;

\textup{II}\,$(V, W)$;

\textup{III}\,$(V, W)$ : le support de $V$ est une variété d'appui $L$ de $W$, et $V$ est
l'intersection de $W$ avec un demi-espace $H$ de $L$ (éventuellement, $H\cap W$ peut être égal à
$L\cap W$).
\end{theoreme}
% (No proof of Theorem 1 is printed beyond « Des propositions 7, 8 et 4 résulte le ».)

Dans le cas projectif, on a le

% [sic] « THÉORÈME 1. » : this is the théorème 1′ announced at the head of § 5 (« les théorèmes
% 1 et 1′ »); the printed number is kept as the argument, the label says 1prime.
\begin{theoreme}[1]
Soit $W$ un corps convexe de $S(n)$, $V$ une partie compacte convexe non vide de $W$; alors les
trois propriétés ci-dessous sont équivalentes :

\textup{I}$'$\,$(V, W)$ : $W-V$ est convexe;

\textup{II}$'$\,$(V, W)$ : les hémisphères ouverts contenant $V$ coupent $W$ suivant un système
fondamental de voisinages de $V$ considéré comme partie de $W$;

\textup{III}$'$\,$(V, W)$ : le support de $V$ est une variété d'appui $L$ de $W$, et $V$ est
l'intersection de $W$ avec un hémisphère $H$ de $L$.
\end{theoreme}

% --- p. 281 ---
\begin{proof}
En effet, le corps convexe $W$ appartient à un hémisphère ouvert de $S(n)$ muni canoniquement d'une
structure d'espace affine à $n$ dimensions. Dans cet espace, \bzcf{le théorème
1$'$}{\ref{bzthm:c3s5t1prime_reunion_convexes_sphere}} est un corollaire du \bzcf{théorème
1}{\ref{bzthm:c3s5t1_reunion_convexes_affine}}.
\end{proof}
% (« le théorème 1′ » names the statement this paragraph proves: it is kept marked, because the
% printed heading of that statement says « 1 », see the [sic] above.)

\subsection*{6. Prolongement d'un convexe au-delà d'un hémisphère.}

\begin{defin}
Soit $W$ un convexe contenu dans un hémisphère fermé $H$ de $S(E)$; on dit qu'un convexe $V$
\textit{prolonge} $W$ au-delà de $\mathring{H}$ si si $V\cap\mathring{H}=W\cap\mathring{H}$. Si
$\mathring{H}\cap W\neq\emptyset$ et si $V$ et $W$ sont fermés, on a de plus :
$W\cap H=V\cap H$.
% [sic] « si si » : the word is printed twice.
\end{defin}

\begin{prop}[1]
Soit $W$ un convexe fermé de noyau vide (\bzcf{\textit{cf.} \S\,2,
déf. 4}{\ref{bzdef:c3s2d4_conjugues}}), et soit $H$ un hémisphère fermé de frontière $h$, tel que
$W\subset H$, $W\cap\mathring{H}\neq\emptyset$; alors il existe un plus grand convexe $V$
prolongeant $W$ au-delà de $\mathring{H}$; $V$ est appelé \textit{saturé de $W$ au-delà} de
$\mathring{H}$. Le noyau de $V$ est vide, sauf si $W$ est l'enveloppe convexe d'un point de
$\mathring{H}$ et de $W\cap h$; ce point de $\mathring{H}$ est alors unique.
\end{prop}
% [sic] « cf. § 2, déf. 4 » : the noyau is defined in § 1, déf. 5 (bzdef:c3s1d5_noyau); § 2,
% déf. 4 is « L et L′ conjugués ». Following the convention, the \ref names the PRINTED § 2,
% déf. 4; see NOTES.md (to be ruled).

\begin{proof}
\textsc{Démonstration} :

1\textsuperscript{o} La famille des convexes prolongeant $W$ au-delà de $\mathring{H}$ est
filtrante croissante.

En effet, soient $V'$, $V''$ deux convexes prolongeant $W$ au-delà de $\mathring{H}$;
l'enveloppe convexe $V'''$ de $V'$ et $V''$ prolonge $W$ au-delà de $\mathring{H}$ [$V'''$ est
réunion de $V'$, de $V''$ et de l'ensemble des segments $(v', v'')$ joignant un point $v'$ de $V'$
à un point $v''$ de $V''$ non diamétralement opposé à $v'$; or
$(v', v'')\cap\mathring{H}\subset W$ : en effet, si $(v', v'')$ rencontre $\mathring{H}$, on a
$v'\in\mathring{H}$ ou $v''\in\mathring{H}$; dans le premier cas, par exemple, on a aussi
$v'\in V''$, donc
\[ (v', v'')\subset V''\qquad\text{et}\qquad(v', v'')\cap\mathring{H}\subset
V''\cap\mathring{H}=W\cap\mathring{H}]. \]

2\textsuperscript{o} Si $V'$ prolonge $W$ au-delà de $\mathring{H}$, il en est de même de
$\overline{V'}$.

Cela résulte de ce que $\overline{V'}$ est convexe et du fait que
$\overline{V'}\cap\mathring{H}$ est l'adhérence, dans $\mathring{H}$, de
$V'\cap\mathring{H}=W\cap\mathring{H}$, qui est fermé dans $\mathring{H}$.

3\textsuperscript{o} Il existe un convexe maximal $V$ prolongeant $W$ au-delà de $\mathring{H}$,
et $V$ est fermé.

Pour construire $V$, ordonnons par inclusion l'ensemble des convexes prolongeant $W$ au-delà de
$\mathring{H}$. L'ensemble ordonné ainsi obtenu est inductif parce que la réunion d'une famille
croissante de convexes prolongeant $W$ au-delà de $\mathring{H}$ est un convexe prolongeant $W$
au-dela de $\mathring{H}$.
% [sic] « au-dela » : read « au-delà ».

Qu'un convexe $V$ maximal parmi ceux qui prolongent $W$ au-delà de $\mathring{H}$ soit le plus
grand convexe jouissant de la proprieté et qu'il soit fermé résulte de 1\textsuperscript{o} et
2\textsuperscript{o}.
% [sic] « proprieté » : read « propriété ».

% --- p. 282 ---
4\textsuperscript{o} Si $V$ contient un couple de points diamétralement opposés $v$ et $v'$, l'un
de ces deux points, $v$, appartient à $\mathring{H}$, et $W$ est l'enveloppe convexe de $v$ et de
$W\cap h$.

En effet, il est impossible que $v$ et $v'$ appartiennent à $h$, car autrement, si $w$ désigne un
point de $W\cap\mathring{H}$, on aurait
\[ \begin{gathered} )w, v(\subset W,\\ )w, v'(\subset W, \end{gathered} \]
d'où, puisque $W$ est fermé,
\[ \begin{gathered} v\in W,\\ v'\in W, \end{gathered} \]
ce qui est contraire à l'hypothèse.

Si $v\in\mathring{H}$, $V$ est réunion de demi-droites fermées joignant $v$ et $v'$, et $W$ n'est
donc autre que le cône de sommet $v$ et de base $W\cap h$. Le point $v$ est unique; en effet,
supposons qu'il existe, dans $\mathring{H}$, $u\neq v$ tel que $W$ soit le cône de sommet $u$ et
de base $W\cap h$; la droite joignant $u$ et $v$ couperait $h$ en deux points diamétralement
opposés qui devraient être tous deux dans $W\cap h$, ce qui est contraire à l'hypothèse faite que
le noyau de $W$ est vide.
\end{proof}

% « COROLAIRE. » : the heading is printed with one L ([sic], read « COROLLAIRE »).
\begin{corollaire}
Soit $W$ un corps convexe de $S(n)$, $H$ un hémisphère fermé de frontière $h$ contenant $W$; si
$W$ n'est pas l'enveloppe convexe d'un point et de $W\cap h$, le saturé de $W$ au-delà de
$\mathring{H}$ est un corps convexe.
\end{corollaire}
% (No proof of the corollary is printed.)

\begin{prop}[2]
Soit $W$ un corps convexe de $S(n)$, $H'_i$ une suite d'hémisphères fermés telle que $H'_i\cap W$
tende vers un point frontière $f$ de $W$. Posons $H_i=\complement\mathring{H}'_i$; alors le saturé
$V_i$ de $W\cap H_i$ au-delà de $\mathring{H}_i$ tend vers $W$.
\end{prop}

\begin{proof}
Supposons qu'il n'en soit pas ainsi; on pourrait trouver une sous-suite $\{V_j\}$ de la suite
$\{V_i\}$ et dans chaque $V_j$ un point $v_j$ de sorte que $v_j$ tende vers $v$ n'appartenant pas
à $W$, et que $H_j$ tende vers $H$ de frontière $h$. Nous allons montrer que cela est absurde. Soit
$h_j$ la variété frontière de $H_j$. L'enveloppe convexe de $v_j$ et de $W\cap H_j$ coupe $h_j$
suivant un ensemble $F_j$ qui est inclus dans
\[ V_j\cap h_j=W\cap h_j; \]
$F_j$ tend vers l'intersection de $h$ et de l'enveloppe convexe de $v$ et de $W$; cependant,
$W\cap h_j$ tend vers le point $f$, d'où la contradiction annoncée, car l'enveloppe convexe de $v$
et de $W$ contient un ouvert dont l'intersection avec $h$ ne peut se réduire à un seul point.
\end{proof}

\begin{prop}[3]
Soit $W_1$ (resp. $W_2$) un corps convexe de $S(n)$; $H_1$ (resp. $H_2$) un hémisphère fermé tel
que $W_1\subset H_1$ (resp. $W_2\subset H_2$) et que le saturé $V_1$ (resp. $V_2$) de $W_1$ (resp.
$W_2$) au-delà de $\mathring{H}_1$ (resp. $\mathring{H}_2$) soit
% --- p. 283 ---
distinct de $W_1$ (resp. $W_2$). On peut trouver un corps convexe $W'$ et deux hémisphères fermés
$H'_1$, $H'_2$ tels que :

1\textsuperscript{o}
\[ \begin{aligned} &H'_1\cap W'=W'_1;\\ &H'_2\cap W'=W'_2,\qquad W'_1\cap W'_2=\emptyset,
\end{aligned} \]
$W'$ étant l'enveloppe convexe de $W'_1$ et de $W'_2$.

2\textsuperscript{o} Il existe deux transformations projectives $g_1$, $g_2$ de l'espace telles
que
\[ \begin{aligned} &g_1H_1=H'_1;\qquad g_1W_1=W'_1;\\ &g_2H_2=H'_2;\qquad g_2W_2=W'_2.
\end{aligned} \]
Nous dirons que $W'$ est obtenu en \textit{raccordant} $W_1$ et $W_2$ suivant les intersections
de ces corps avec les hyperplans $h_1$, $h_2$ frontières de $H_1$, $H_2$; si $g_1$ est
l'identité, nous dirons que $W'$ est obtenu en \textit{accolant} $W_2$ à $W_1$ suivant
$h_1\cap W_1$, $h_2\cap W_2$.
\end{prop}

\begin{proof}
\textsc{Démonstration.} --- Soit $a''_1$ un point de $V_1$, $a''_1\notin W_1$; soit $a_1$ le point
diamétralement opposé à $a''_1$; $W_1$ est contenu dans l'enveloppe convexe $C_1$ du point $a_1$
et de $F_1=h_1\cap W_1$. En effet, soit $m$ un point de $W_1$, le segment $(a''_1\,m)$ coupe $F_1$
en un point $f$ et $m$ appartient au segment $(a_1\,f)$.

On pose $F_2=W_2\cap h_2$ et l'on choisit de même $a_2$ en sorte que l'enveloppe convexe $C_2$ de
$a_2$ et de $F_2$ contienne $W_2$.

Nous allons raccorder $C_1$ et $C_2$ suivant $F_1$, $F_2$. Il est clair que les transformations
$g_1$, $g_2$ utilisées permettront de raccorder $W_1$ et $W_2$, d'où la proposition.

Faisons choix d'un hémisphère ouvert $\mathring{H}$ muni canoniquement d'une structure d'espace
affine. Soit $h'_1$, $h'_2$ deux hyperplans parallèles; $F'_1$ (resp. $F'_2$) un convexe de $h'_1$
(resp. $h'_2$) projectivement équivalent à $F_1$ (resp. $F_2$). Choisissons $a'_1\notin h'_1$
(resp. $a'_2\notin h'_2$) de sorte que la réunion de toutes les demi-droites issues de $a'_1$
(resp. $a'_2$) et rencontrant $F'_1$ (resp. $F'_2$) contiennent $F'_2$ (resp. $F'_1$). L'enveloppe
convexe $C'$ de $a'_1$, $a'_2$, $F'_1$, $F'_2$, raccorde $C_1$ et $C_2$ suivant $F_1$ et $F_2$.
\end{proof}

\begin{prop}[4]
Sur tout corps convexe $W$ de $S(n)$ ($n\geqq 1$), on peut trouver deux facettes distinctes dont
chacune est réduite à un point.
\end{prop}

\begin{proof}
Remarquons que la notion de facette réduite à un point s'identifie à celle de point extrémal. La
proposition va résulter du théorème de Krein-Milman (\textit{cf.} \textsc{Bourbaki},
\textit{Espaces vectoriels topologiques}, chap. II, \S\,4, théorème 1) dont nous rappelons
l'énoncé : « Soit $E$ un espace vectoriel topologique localement convexe séparé; tout convexe
compact de $E$ est l'enveloppe convexe fermée de l'ensemble de ses points extrémaux ».

Soit $\mathring{H}$ un hémisphère ouvert de $S(n)$ contenant $W$, $\mathring{H}$ est muni
canoniquement d'une structure d'espace affin, il peut être muni d'une structure d'espace
vectoriel à $n$ dimensions $E$. Il suffit d'appliquer au convexe compact $W$ de $E$ le théorème de
Krein-Milman.
\end{proof}

% (Chapter 3 ends at the foot of p. 283. Chapter 4 « Corps et formes convexes affines » begins at
% the top of p. 284; not transcribed in this file.)

\end{document}
